\section{The baseline, and the shape of the failure}
\label{sec:baseline}

\begin{table}[h]
\centering
\begin{tabular}{lrr}
\toprule
& \textbf{single-hop (\rQsingle{}\,Q)} & \textbf{multi-hop (\rQmulti{}\,Q)} \\
\midrule
recall@5\ \ \textbf{all}   & \rSallFive{}\%   & \rMallFive{}\% \\
recall@10\ \textbf{all}    & \rSallTen{}\%    & \rMallTen{}\% \\
recall@20\ \textbf{all}    & \textbf{\rSallTwenty{}\%} & \textbf{\rMallTwenty{}\%} \\
\midrule
recall@5\ \ \textbf{any}   & \rSanyFive{}\%   & \rManyFive{}\% \\
recall@10\ \textbf{any}    & \rSanyTen{}\%    & \rManyTen{}\% \\
recall@20\ \textbf{any}    & \rSanyTwenty{}\% & \rManyTwenty{}\% \\
\bottomrule
\end{tabular}
\caption{Dense retrieval on LoCoMo, \texttt{zen-embedding-0.6b}, brute-force
cosine \meas{}. The bottom row is the one that reframes the problem.}
\label{tab:baseline}
\end{table}

Read the table by column and it says the retriever is bad at multi-hop
questions: \rSallTwenty{}\% against \rMallTwenty{}\% at $k{=}20$ is a collapse of
more than three to one.

Read it by row and it says something different and more useful. On the
\textbf{any} rows the two columns converge: \rSanyTwenty{}\% and
\rManyTwenty{}\% at $k{=}20$ are within a point of each other. Whatever the
retriever is doing, it does it equally well on both categories --- it puts a
relevant turn in the top 20 about four times in five, regardless of how many
turns the question actually needs.

So the failure is not comprehension and not embedding quality. Those would
depress \textbf{any} as well, and \textbf{any} is not depressed. The failure is
specifically in \emph{completeness}: the gap between \rManyTwenty{}\% and
\rMallTwenty{}\% --- some 57 percentage points --- is questions where the
retriever found part of the evidence and missed the remainder.

\paragraph{Why more $k$ is not obviously the answer.}
The single-hop column improves steadily with $k$ (\rSallFive{} $\to$
\rSallTen{} $\to$ \rSallTwenty{}), which is what a retriever that is
merely under-budgeted looks like. The multi-hop \textbf{all} column improves too
(\rMallFive{} $\to$ \rMallTen{} $\to$ \rMallTwenty{}) but from so low a base that
extrapolating it to a useful number requires a $k$ that defeats the purpose of
retrieving. \S\ref{sec:diagnostic} measures how large that $k$ would have to be,
and the answer rules the approach out.
